\chapter{Experiment Parameters and Runtimes}
\label{appn:params}

%% -----------------------------------------------------------------------
\section*{E0 --- Analytical Baseline}

E0 uses newsvendor base-stock levels computed analytically from demand
statistics and lead times.  No dispatch threshold is set (all edges use
\texttt{min\_dispatch\_utilization = 0.0}).

%% -----------------------------------------------------------------------
\section*{E3 --- Joint Optimisation}

Base-stock levels (integers) and dispatch thresholds (fractions)
optimised by Optuna TPE across 100 trials with a 92\% fill-rate
constraint.  Best feasible trial: \#62.

\subsection*{Base-stock levels (per node, per SKU)}

\begin{table}[H]
  \centering
  \caption{E3 optimised base-stock levels.}
  \begin{tabular}{@{}lccccc@{}}
    \toprule
    Node & SKU\_A & SKU\_B & SKU\_C & SKU\_D & SKU\_E \\
    \midrule
    W1 & 1639 & 1022 & 1172 &  595 & 1563 \\
    W2 & 3703 &  769 & 2005 &  832 & 1279 \\
    R1 &  929 &  567 & 1028 &  504 &  249 \\
    R2 & 1226 &  781 &  757 &  442 &  584 \\
    R3 & 1620 &  653 & 1607 &  939 &  979 \\
    \bottomrule
  \end{tabular}
\end{table}
\vspace{-1em}
\subsection*{Dispatch thresholds (per lane)}

\begin{table}[H]
  \centering
  \caption{E3 optimised dispatch thresholds.}
  \begin{tabular}{@{}lc@{}}
    \toprule
    Lane & \texttt{min\_dispatch\_utilization} \\
    \midrule
    Supplier $\to$ W1 & 0.792 \\
    Supplier $\to$ W2 & 0.361 \\
    W1 $\to$ R1       & 0.635 \\
    W1 $\to$ R2       & 0.470 \\
    W2 $\to$ R3       & 0.818 \\
    \bottomrule
  \end{tabular}
\end{table}

The long-haul Supplier$\to$W1 lane and the W2$\to$R3 lane carry the
highest thresholds ($\sim$0.79--0.82), consolidating nearly full trucks.
The Supplier$\to$W2 lane operates at a lower threshold (0.36), reflecting
the smaller demand served by W2.  Warehouse-to-retailer thresholds
(0.47--0.82) balance consolidation against service frequency.

%% -----------------------------------------------------------------------
\section*{E9 --- Stochastic Demand Robustness}

E9 holds the E3 parameters fixed and evaluates them under Poisson demand,
replacing each M5 demand series with a Poisson generator whose rate
$\lambda$ equals the mean daily demand over the 365-day evaluation window.
Ten independently seeded realisations are run; no re-optimisation is
performed.

\begin{table}[H]
  \centering
  \caption{E9 stochastic robustness results (10 Poisson seeds, E3 params).}
  \begin{tabular}{@{}lcc@{}}
    \toprule
    Metric & Mean & Std \\
    \midrule
    Total cost (\$) & 242{,}000 & 1{,}000 \\
    Fill rate (\%)  & 99.95     & 0.03    \\
    \bottomrule
  \end{tabular}
\end{table}

The 92\% fill-rate target is met in all ten realisations, confirming that
the deterministically-optimised E3 parameters transfer to stochastic demand
without re-optimisation.

%% -----------------------------------------------------------------------

\section*{Empirical Runtimes}
\label{appn:runtimes}

All timings were measured on an Intel Core i7-14700HX (20 cores,
single socket) running Python~3.12 with Optuna~3.x in single-threaded
mode.  Wall-clock time was captured with \texttt{time.perf\_counter()}
around each optimiser or simulation call; file I/O and plotting are
excluded.

\subsection*{Per-trial cost by optimisation mode}

\begin{table}[H]
  \centering
  \caption{Per-trial wall-clock time (20-trial sample, averaged).}
  \label{tab:rt-per-trial}
  \begin{tabular}{@{}lcr@{}}
    \toprule
    Mode & Dimensions & ms\,/\,trial \\
    \midrule
    Transport-only & 5  & 127 \\
    Inventory-only & 25 & 140 \\
    Joint          & 30 & 131 \\
    \bottomrule
  \end{tabular}
\end{table}
Per-trial cost is flat from 5 to 30 dimensions (127--140\,ms),
confirming that TPE bookkeeping overhead is negligible relative to the
365-day simulation evaluation ($\approx130$\,ms per run).

\subsection*{Per-experiment runtimes at the canonical trial budget}

\begin{table}[H]
  \centering
  \caption{Per-experiment wall-clock time (\texttt{--trials~100}, single-threaded).}
  \label{tab:rt-per-exp}
  \begin{tabular}{@{}llcrr@{}}
    \toprule
    Experiment & Mode & Dims & Sim.\ evals & Wall-clock (s) \\
    \midrule
    E0  & Baseline (1 sim)                 & ---    &    1 &  $<1$ \\
    E1a & Fixed threshold (1 sim)          & ---    &    1 &  $<1$ \\
    E1b & Transport-only                   & 5      &  100 &   13  \\
    E2  & Inventory-only                   & 25     &  100 &   14  \\
    E3  & Joint                            & 30     &  100 &   13  \\
    E3\_per\_sku & Joint (per-SKU)         & 30     &  100 &   13  \\
    E4  & NSGA-II + sweep                  & 30     & 1000 &  130  \\
    E5  & Disruption (2 sims)              & ---    &    2 &  $<1$ \\
    E6  & Policy search (4$\times$100)     & 30--35 &  400 &   53  \\
    E7  & Forecast sweep (5 sims)          & ---    &    5 &    1  \\
    E8  & Bullwhip-aware (5$\times$100)    & 30     &  500 &   65  \\
    E9  & Stochastic (10 sims)             & ---    &   10 &    1  \\
    \midrule
    \textbf{Full suite} & & & \textbf{$\sim$2220} & \textbf{$\sim$305} \\
    \bottomrule
  \end{tabular}
\end{table}

The full E0--E9 suite completes in approximately \textbf{5 minutes} on
a laptop CPU.  Peak heap allocation per evaluation is \textbf{11.4\,MB};
total process memory stays below 200\,MB throughout the full suite.
